
\documentclass[11pt,a4paper]{article}
\usepackage[utf8]{inputenc}
\usepackage[T1]{fontenc}
\usepackage{lmodern}
\usepackage[margin=2.2cm]{geometry}
\usepackage{microtype}
\usepackage{graphicx}
\usepackage{booktabs}
\usepackage{longtable}
\usepackage{tabularx}
\usepackage{array}
\usepackage{float}
\usepackage{fancyhdr}
\usepackage{xcolor}
\usepackage{hyperref}
\usepackage{enumitem}
\usepackage{setspace}
\usepackage{caption}
\usepackage{titlesec}
\usepackage{amsmath}
\usepackage{pdflscape}
\usepackage{verbatim}
\definecolor{accent}{HTML}{1F4E79}
\definecolor{soft}{HTML}{EAF1F7}
\definecolor{textgray}{HTML}{4D5966}
\hypersetup{colorlinks=true, linkcolor=accent, citecolor=accent, urlcolor=accent, pdftitle={MAS LCE project compendium}}
\pagestyle{fancy}
\fancyhf{}
\fancyfoot[C]{\thepage}
\renewcommand{\headrulewidth}{0pt}
\setstretch{1.06}
\titleformat{\section}{\Large\bfseries\color{accent}}{\thesection}{0.6em}{}
\titleformat{\subsection}{\large\bfseries\color{accent}}{\thesubsection}{0.6em}{}
\captionsetup{font=small,labelfont=bf}
\setlength{\parskip}{0.5em}
\setlength{\parindent}{0pt}

\begin{document}
{\color{accent}\rule{\textwidth}{1.2pt}}
\begin{center}
    {\LARGE\bfseries Project guide}\\[0.4em]
    {\large Python-only reproduction package for the MAS LCE fraudulent-classified-ads study}\\[0.8em]
    {\normalsize March 2026}
\end{center}
{\color{accent}\rule{\textwidth}{0.6pt}}

\section*{What this package is}
This package is a self-contained English compendium built from the archived David Graz MAS LCE project. It contains:
\begin{itemize}[leftmargin=1.5em]
    \item an article-style reconstruction of the thesis;
    \item a dedicated integrity and reproducibility audit;
    \item sanitised canonical datasets that are sufficient to reproduce the published descriptive tables and the archived predictive metrics;
    \item Python scripts that regenerate tables, figures, manifests, checksums, and PDF reports;
    \item metadata files that document provenance, version drift, and canonicalisation choices.
\end{itemize}

\section{Package structure}
\begin{tabularx}{\textwidth}{p{4.0cm} X}
\toprule
Path & Contents \\
\midrule
reports/ & Final PDFs and LaTeX sources for the article, audit, and guide \\
figures/ & Reproduced figures used by the reports \\
data/canonical/ & Sanitised canonical exploratory and case-study datasets \\
data/tables/ & CSV tables used in the reports \\
data/json/ & JSON summaries used to populate the LaTeX reports \\
data/metadata/ & Source inventory, discrepancy logs, checksums, manifest, and data dictionary \\
scripts/ & Python-only build and analysis scripts \\
\bottomrule
\end{tabularx}

\section{How to rebuild the package}
The package is designed so that the main build command is a Python script.

\begin{verbatim}
python scripts/build_all.py --use-packaged-data
\end{verbatim}

If you also have access to the original extracted \texttt{\_Master\_LCE\_Docs} directory, you can regenerate the canonical datasets and the inventory directly from the source corpus:

\begin{verbatim}
python scripts/build_all.py --source-root /path/to/_Master_LCE_Docs
\end{verbatim}

The build script regenerates:
\begin{itemize}[leftmargin=1.5em]
    \item canonical CSV datasets;
    \item descriptive and discrepancy tables;
    \item ROC and indicator-distribution figures;
    \item LaTeX files and compiled PDFs;
    \item a manifest and SHA-256 checksums for all package files.
\end{itemize}

\section{Canonical datasets}
Two canonical datasets are bundled because the archive contains multiple versions of the same analyses.
\begin{itemize}[leftmargin=1.5em]
    \item \texttt{exploratory\_sample\_canonical.csv} reproduces the exploratory verification sample used by the thesis and by the archived exploratory random-forest script.
    \item \texttt{case\_study\_sample\_canonical.csv} reproduces the broader case-study sample whose counts and predictive metrics match the final thesis narrative.
\end{itemize}

Both canonical files are sanitised: personally identifying profile fields, links, message fragments, and contact artefacts were removed because they are not necessary for the quantitative reproduction.

\section{Data dictionary}
\begin{tabularx}{\textwidth}{p{4.0cm} X}
\toprule
Field & Meaning \\
\midrule
listing\_id & Stable listing identifier inside the archived dataset \\
source\_status & Archived manual status (POSITIVE, NEGATIVE, IDENTIFIED, DETECTED) \\
ad\_text & Listing description text \\
price\_chf & Listing price in CHF-equivalent numeric form \\
profile\_creation\_year & Year the seller profile was created; missing for opaque profiles \\
content\_score / price\_score / profile\_score & Encoded operational indicators used by the original project \\
fraud\_index & Sum of indicator scores used by the thesis \\
risk\_target & Binary operational target: 0 = acceptable, 1 = unfavorable \\
price\_ratio & Derived ratio between listing price and the model-specific market benchmark \\
\bottomrule
\end{tabularx}

\section{Integrity and ethics notes}
The source archive contains 1024 files, including legacy materials that should not be reused operationally. In particular:
\begin{itemize}[leftmargin=1.5em]
    \item legacy scraper scripts with embedded credentials were excluded from the compendium;
    \item early similarity scripts with self-comparison bugs were excluded from the final analysis pipeline;
    \item reproduction is intentionally offline and depends on archived snapshots rather than fresh scraping.
\end{itemize}

\section{Reading order}
For most users, the recommended order is:
\begin{enumerate}[leftmargin=1.5em]
    \item read the article-style reconstruction;
    \item read the integrity audit to understand what was reproduced and what remains uncertain;
    \item inspect the canonical datasets and the metadata tables if you need technical traceability;
    \item run \texttt{python scripts/build\_all.py --use-packaged-data} if you want to regenerate all outputs.
\end{enumerate}

\end{document}
